\honest{Why Truth?}{Eliezer Yudkowsky}{2006}

Some commenters have asked why rationalists care about truth. I give three answers, and say that each one shapes the search differently.

First, you might think seeking truth is noble. Then you will probably expect others to look behind the curtain too. I am suspicious of this motive. I do not reject the ideal, but it is too easy to learn, as moral duties, ``modes of thinking that are dreadful missteps in the dance.'' My example is Spock, who is always calm and gives probabilities with too many digits. In a footnote he gives the Enterprise a 2.234\% chance of surviving, and the ship survives nine times out of ten. Many people think this is what the duty to be rational looks like. \nb{Spock shows a mistaken idea of rationality. He does not show that treating rationality as a duty produces it.} A moral duty has ``all the dreadful degrees of freedom of an arbitrary tribal custom,'' and people who get the wrong answer protest that they acted with propriety. I give no case of this.

Second, you might want truth for a goal: an airplane, or chocolate milk. Then you rank questions by the expected value of the information: how much the possible answers affect your choices, how much the choices matter, and how likely you are to find an answer that changes what you would otherwise do. This may seem impure, but it gives an outside check. If the plane falls, or the shop has no chocolate milk, you learn that you thought badly. So you learn which ways of thinking work. Of my three motives, this is the only one I give a test.

Third, curiosity. It is pure, and it ranks questions by how interesting they are. A harder problem with more chance of failure may be worth more effort just because it is more fun. But it ``may not linger too long on verifying its answers.'' What set humanity on the path of science was noticing that some ways of thinking let us ``manipulate the world.'' Campfire tales of gods and heroes satisfied curiosity ``just as well, and no one realized that anything was wrong with that.'' I cite no historian. \nb{If this covers the ancient Greeks, it is overstated. In the late sixth century BC Xenophanes complained about the poets' stories of the gods, and noted that each people pictures its gods in its own likeness.}

To improve beyond the hunter-gatherer standard we need deliberate rules of thinking, ``things that look like norms of rationalist `propriety.'\,'' A new habit of thought starts as an explicit rule and only slowly, if ever, becomes part of how we are moved. \nb{Propriety is what the post blamed a few paragraphs earlier, in people who defend a wrong answer by saying they acted with propriety.}

All three, curiosity, practical use and these rules, belong to the project. Asked which is most foundational, I answer ``curiosity.'' My reason is that ``I also just really want to know.'' \nb{The post's own history gave the credit for science to the practical motive, and said that campfire tales satisfied curiosity.}
